%============================================================
% CH2 · PATTERN 5 — REVERSE SEARCH: THEORY
%============================================================
\subsection{Pattern Theory}
\label{ssec:ch2-p5-theory}

% local styles of this pattern: a reversed edge, and the whole graph D drawn reversed
\tikzset{
    g2p5-re/.style={ge, {Stealth[length=2.4mm]}-},
    g2p5-rt/.style={gtree, {Stealth[length=2.4mm]}-},
    g2p5-Drev/.style={g2-e12/.style=g2p5-re, g2-e15/.style=g2p5-re, g2-e16/.style=g2p5-re, g2-e23/.style=g2p5-re,
        g2-e24/.style=g2p5-re, g2-e31/.style=g2p5-re, g2-e43/.style=g2p5-re, g2-e54/.style=g2p5-re,
        g2-e56/.style=g2p5-re, g2-e67/.style=g2p5-re, g2-e76/.style=g2p5-re, g2-e87/.style=g2p5-re},
    % grid cells, used as  \foreach \s [count=\c from 0] in {w,e,o,...} \node[g2p5-\s] at (\c,-\r) {};
    g2p5-e/.style={gcell, minimum size=6.5mm, font=\scriptsize\bfseries},
    g2p5-w/.style={g2p5-e, fill=midgray!70},
    g2p5-o/.style={g2p5-e, fill=computedfill},
    g2p5-i/.style={g2p5-e, fill=answerfill},
    g2p5-G/.style={g2p5-e, fill=basefill, label={center:G}},
    g2p5-B/.style={g2p5-e, fill=dupfill, label={center:B}},
    g2p5-Go/.style={g2p5-e, fill=computedfill, draw=successgreen, line width=1.1pt, label={center:G}},
    g2p5-x/.style={g2p5-e, fill=midgray!70, draw=recursionred, line width=1.1pt}
}

\begin{patternbox}{Pattern 5 at a Glance}
``Which vertices can reach the target $t$?'' is $n$ questions, one per start vertex. Turn every edge around: $u$ reaches $t$ in $G$ \textbf{iff} $t$ reaches $u$ in the reversed graph $G^R$. One search from $t$ in $G^R$ answers all $n$ questions at once --- and a search from a \emph{set} of targets answers ``who can reach any of them?'' just as cheaply.

\begin{center}
\begin{tikzpicture}[x=1.0cm, y=0.95cm]
    \begin{scope}
        \pic {g2-D};
        \node[font=\scriptsize\bfseries] at (2.4,2.25) {$D$: who can reach 4? ($8$ forward searches)};
    \end{scope}
    \begin{scope}[xshift=7.2cm, g2p5-Drev, g2-v4/.style=gvdst, g2-v1/.style=gvdone, g2-v2/.style=gvdone,
                  g2-v3/.style=gvdone, g2-v5/.style=gvdone, g2-v6/.style=gvdead, g2-v7/.style=gvdead, g2-v8/.style=gvdead,
                  g2-e24/.style=g2p5-rt, g2-e54/.style=g2p5-rt, g2-e15/.style=g2p5-rt, g2-e31/.style=g2p5-rt]
        \pic {g2-D};
        \node[font=\scriptsize\bfseries] at (2.4,2.25) {$D^R$: one search from 4};
    \end{scope}
    \draw[flow] (5.3,0.2) -- node[lbl, above]{reverse} (6.8,0.2);
\end{tikzpicture}
\end{center}
\begin{recurrencebox}
\centering\small
build \texttt{radj[v]} $= \{u : u \to v\}$\quad;\quad push every target, mark it\quad;\quad pop $v$, push every unmarked $u \in \texttt{radj}[v]$\\[3pt]
marked $=$ $\{u : u \leadsto T\}$ \qquad one search, \VE{} instead of $n$ searches, \bigO{n(n+m)}
\end{recurrencebox}
\end{patternbox}

%------------------------------------------------------------
\subsubsection*{From the lessons to the pattern}

Tool~5 of Lesson~\ref{ssec:ch2-toolbox} showed the trick in two lines. This pattern is about \emph{seeing} it in a statement, because the statement almost never says ``reversed graph''. It says ``every good person must be able to escape'', ``water flows to the ocean'', ``a region is captured unless it touches the border'', ``list all ancestors''. In each of these the question is asked from the point of view of many starting points, all heading to the same few places. The pattern is to change the point of view:
\begin{itemize}[leftmargin=*]
    \item \textbf{reverse the edges} when the graph is directed (an edge $u \to v$ becomes $v \to u$; in a grid with rules like ``water flows downhill'' the reversed rule is ``climb uphill'');
    \item \textbf{reverse the roles} when the graph is undirected: the search does not have to change at all, only its start --- from the exit, the border, the ocean, instead of from every cell;
    \item \textbf{start from all targets at once} (several exits, the whole border, every cell of a set $T$): one search with many sources costs the same as a search with one.
\end{itemize}

\begin{diagrambox}
\begin{tikzpicture}[x=1cm, y=1cm]
    \begin{scope}
        \node[font=\small\bfseries] at (2.0,1.55) {forward: one search per start};
        \foreach \i/\y in {1/1.0, 2/0.35, 3/-0.3, 4/-0.95}{
            \node[gvsrc, minimum size=6mm, font=\scriptsize\bfseries] (s\i) at (0,\y) {$s_\i$};
            \draw[gtree, -{Stealth[length=2.2mm]}, decorate, decoration={snake, amplitude=0.6mm, segment length=4mm, post length=2mm}] (s\i) -- (3.4,0.05);
        }
        \node[gvdst] (t) at (3.8,0.05) {$t$};
        \node[note, text width=4.2cm] at (2.0,-1.65) {$n$ searches, each may cost \bigO{n+m}: \bigO{n(n+m)}};
    \end{scope}
    \begin{scope}[xshift=7.6cm]
        \node[font=\small\bfseries] at (2.0,1.55) {reverse: one search from $t$};
        \foreach \i/\y in {1/1.0, 2/0.35, 3/-0.3, 4/-0.95}{
            \node[gvdone, minimum size=6mm, font=\scriptsize\bfseries] (r\i) at (0,\y) {$s_\i$};
        }
        \node[gvdst] (t2) at (3.8,0.05) {$t$};
        \foreach \i/\y in {1/1.0, 2/0.35, 3/-0.3, 4/-0.95}
            \draw[gtree, {Stealth[length=2.2mm]}-] (r\i) -- (t2);
        \node[note, text width=4.2cm] at (2.0,-1.65) {one search in $G^R$, every vertex and edge once: \VE};
    \end{scope}
\end{tikzpicture}
\figcap{The same question asked from both ends. The answer is the same set of vertices; only the cost differs.}
\end{diagrambox}

%------------------------------------------------------------
\subsubsection*{When to recognise it (and when not)}

\begin{itemize}[leftmargin=*]
    \item ``\textbf{from which} vertices / cells can you reach the exit, the border, the ocean, a vertex of $T$?'' --- one search from the targets in $G^R$;
    \item ``\textbf{for every} vertex, all its ancestors / all vertices that can reach it'' --- one reverse search \emph{per vertex}, \bigO{n(n+m)}, fine for $n \le 1000$--$2000$ (Problem~\ref{ssec:ch2-all-ancestors});
    \item ``every good one must get out, no bad one may'' --- first make the bad ones harmless, then one search from the exit checks all good ones (Problem~\ref{ssec:ch2-solve-maze});
    \item ``regions that do \emph{not} touch the border'', ``lakes'', ``captured'', ``surrounded'' --- search from the border; whatever stays unmarked is enclosed;
    \item ``can the walk go on forever?'', ``which states are safe / doomed?'' --- peel dead ends backwards (reverse peeling, below);
    \item ``does everybody reach 1 \emph{and} does 1 reach everybody?'' --- one search from 1 in $G$, one in $G^R$.
\end{itemize}
\textbf{Not this pattern:} \emph{shortest distances} to a target from all vertices need BFS in $G^R$ (and from many sources a multi-source BFS, Chapter~3) or Dijkstra with weights; DFS gives the right \emph{set} but not the distances. Full \emph{strong connectivity} --- which pairs reach each other --- is the strongly connected components of Chapter~19; the two-search test above only answers ``is the whole graph one component?''. Reachability for all $n^2$ pairs on large graphs needs a DAG of components and bitsets, not $n$ searches.

%------------------------------------------------------------
\subsubsection*{The algorithm and its invariant}

\begin{invariantbox}[Why one reverse search is enough]
\textbf{Lemma.} $u \leadsto v$ in $G$ iff $v \leadsto u$ in $G^R$. A path $u = x_0 \to x_1 \to \dots \to x_k = v$ in $G$ is, read backwards, the path $v = x_k \to x_{k-1} \to \dots \to x_0 = u$ in $G^R$, since every edge was turned around; the converse is the same argument applied to $G^R$, whose reverse is $G$.

\textbf{Multi-source.} Add a new vertex $\star$ with edges $\star \to t$ for every $t \in T$ in $G^R$. A search from $\star$ marks exactly the vertices $v$ with $t \leadsto v$ in $G^R$ for some $t$, i.e.\ with $v \leadsto T$ in $G$. Pushing all of $T$ on the stack before the loop is exactly ``entering $\star$''.

\textbf{Invariant of the loop.} Every marked vertex can reach $T$ in $G$ (it was marked along a reversed edge from a marked vertex), and when the stack is empty no reversed edge leaves the marked set --- so, by the closure argument of Chapter~1, every vertex that can reach $T$ is marked.

\textbf{Cost.} Building \texttt{radj} is \bigO{m}; every vertex is pushed once, every reversed edge looked at once: \VE, the same as one forward search.
\end{invariantbox}

\begin{diagrambox}
\begin{tikzpicture}[x=1.0cm, y=0.95cm]
    \begin{scope}[g2p5-Drev, g2-v2/.style=gvdst, g2-v8/.style=gvdst, g2-v1/.style=gvdone, g2-v3/.style=gvdone,
                  g2-v4/.style=gvdone, g2-v5/.style=gvdone, g2-v6/.style=gvdead, g2-v7/.style=gvdead,
                  g2-e12/.style=g2p5-rt, g2-e31/.style=g2p5-rt, g2-e43/.style=g2p5-rt, g2-e54/.style=g2p5-rt]
        \pic {g2-D};
        \node[font=\scriptsize\bfseries] at (2.4,2.25) {$D^R$ from $T = \{2, 8\}$ at once};
    \end{scope}
    \node[note, text width=6.2cm, align=left, anchor=west] at (6.0,0.3) {Stack at the start: \texttt{2 8}. Pop 8: nothing points into 8. Pop 2: mark 1 (edge $1 \to 2$). Pop 1: mark 3 ($3 \to 1$). Pop 3: 2 is marked, mark 4 ($4 \to 3$). Pop 4: 2 marked, mark 5 ($5 \to 4$). Pop 5: 1 is marked.\\[3pt]
    Marked (template output): \texttt{1 2 3 4 5 8}. The targets are marked too --- a vertex of $T$ ``reaches'' $T$ in zero steps. Vertices 6 and 7 cannot reach $T$: their only exits lead to each other.};
\end{tikzpicture}
\figcap{A multi-target reverse search on $D$. Thick blue: the reversed edges along which a new vertex was marked.}
\end{diagrambox}

%------------------------------------------------------------
\subsubsection*{Reversing roles in grids}

In an undirected grid nothing has to be reversed, yet the same idea saves the day: ``which empty cells are \emph{not} connected to the border?'' asked cell by cell is one search per cell, \bigO{(RC)^2}. Asked from the other end, it is one search that starts from \emph{all} empty border cells at once; the unmarked empty cells are the enclosed ones (template 2).

\begin{diagrambox}
\begin{tikzpicture}[x=6.5mm, y=6.5mm]
    \begin{scope}
        \foreach \r/\row in {0/{w,w,w,w,w}, 1/{w,e,e,w,e}, 2/{w,e,w,e,e}, 3/{w,w,w,w,w}}
            \foreach \s [count=\c from 0] in \row \node[g2p5-\s] at (\c,-\r) {};
        \node[font=\scriptsize\bfseries] at (2,1.1) {input (grey $=$ wall)};
    \end{scope}
    \begin{scope}[xshift=5.2cm]
        \foreach \r/\row in {0/{w,w,w,w,w}, 1/{w,i,i,w,o}, 2/{w,i,w,o,o}, 3/{w,w,w,w,w}}
            \foreach \s [count=\c from 0] in \row \node[g2p5-\s] at (\c,-\r) {};
        \draw[flow] (5.1,-1) -- (4.45,-1);
        \draw[flow] (5.1,-2) -- (4.45,-2);
        \node[font=\scriptsize\bfseries] at (2,1.1) {search from the border};
        \node[note, text width=3.6cm, align=left, anchor=west] at (5.9,-1.5) {start: the two empty border cells (arrows); blue $=$ connected to the border, orange $=$ enclosed};
    \end{scope}
\end{tikzpicture}
\figcap{Template 2 on the grid \texttt{\#\#\#\#\# / \#..\#. / \#.\#.. / \#\#\#\#\#}: output \texttt{\#\#\#\#\# \#..\#o \#.\#oo \#\#\#\#\#}. The three orange cells are enclosed. This is Surrounded Regions (E5.1) and the lakes of E5.3.}
\end{diagrambox}

The two-ocean version (E5.2) runs this twice: once from the Pacific border, once from the Atlantic border, with the rule reversed (``flows down'' becomes ``climbs up''), and intersects the two marked sets.

%------------------------------------------------------------
\subsubsection*{Block the bad, then check the good}

Some problems mix two kinds of starts: \emph{good} ones that must reach the target and \emph{bad} ones that must not. Asking ``can good $g$ get out?'' for every $g$ separately is again $n$ searches --- and the answer depends on which cells we decide to block. The reverse view untangles both:
\begin{enumerate}[leftmargin=*]
    \item make every bad start \textbf{harmless} by the cheapest local change that is forced anyway (in Problem~\ref{ssec:ch2-solve-maze}: wall in the four neighbours of every bad person);
    \item then run \textbf{one} search from the target and check that it marks every good start and no bad one.
\end{enumerate}
Step 1 is safe because any solution must already cut every bad start from the target, and a good start whose path passes \emph{next to} a bad one is lost in every solution. The proof is in the problem.

\begin{diagrambox}
\begin{tikzpicture}[x=6.5mm, y=6.5mm]
    \begin{scope}
        \foreach \r/\row in {0/{G,e,e,e}, 1/{e,e,B,e}, 2/{e,w,e,e}}
            \foreach \s [count=\c from 0] in \row \node[g2p5-\s] at (\c,-\r) {};
        \node[font=\scriptsize\bfseries] at (1.5,1.1) {exit $=$ bottom right};
    \end{scope}
    \begin{scope}[xshift=3.3cm]
        \foreach \r/\row in {0/{G,e,x,e}, 1/{e,x,B,x}, 2/{e,w,x,e}}
            \foreach \s [count=\c from 0] in \row \node[g2p5-\s] at (\c,-\r) {};
        \node[font=\scriptsize\bfseries] at (1.5,1.1) {1. wall in B (red)};
    \end{scope}
    \begin{scope}[xshift=6.6cm]
        \foreach \r/\row in {0/{G,e,x,e}, 1/{e,x,B,x}, 2/{e,w,x,o}}
            \foreach \s [count=\c from 0] in \row \node[g2p5-\s] at (\c,-\r) {};
        \node[font=\scriptsize\bfseries] at (1.5,1.1) {2. search from the exit};
    \end{scope}
    \node[note, text width=4.3cm, align=left, anchor=west] at (15.0,-1) {Both neighbours of the exit were walled, so the search from the exit (blue) marks only the exit: G is cut off and the answer is \emph{no} --- every route from G to the exit passes next to B, so no choice of walls does better.};
\end{tikzpicture}
\figcap{The two steps of ``block the bad, check the good''. The blue cells are reached from the exit; G is not among them.}
\end{diagrambox}

%------------------------------------------------------------
\subsubsection*{All ancestors: one reverse search per vertex}

``For every $v$, list all vertices that can reach $v$'' has an output of size up to $n^2/2$, so no algorithm beats $\Theta(n^2)$ in the worst case. One reverse search from every $v$ costs \bigO{n(n+m)} and produces each list already \emph{sorted} if we collect it by scanning $u = 0, 1, \dots, n-1$ (template 3). The mirror image works as well: a forward search from every $u$ in increasing order appends $u$ to every vertex it reaches --- also sorted for free.

\begin{diagrambox}
\begin{tikzpicture}[x=1.15cm, y=1cm]
    \begin{scope}
        \node[gv] (a0) at (0,0.7) {0}; \node[gv] (a1) at (0,-0.7) {1};
        \node[gv] (a2) at (1.4,0.7) {2}; \node[gv] (a3) at (2.8,0.7) {3}; \node[gv] (a4) at (1.4,-0.7) {4};
        \draw[gde] (a0) -- (a2); \draw[gde] (a1) -- (a2); \draw[gde] (a2) -- (a3); \draw[gde] (a1) -- (a4);
        \node[font=\scriptsize\bfseries] at (1.4,1.5) {a DAG};
    \end{scope}
    \begin{scope}[xshift=4.6cm]
        \node[gvdone] (b0) at (0,0.7) {0}; \node[gvdone] (b1) at (0,-0.7) {1};
        \node[gvdone] (b2) at (1.4,0.7) {2}; \node[gvdst] (b3) at (2.8,0.7) {3}; \node[gvdead] (b4) at (1.4,-0.7) {4};
        \draw[g2p5-rt] (b0) -- (b2); \draw[g2p5-rt] (b1) -- (b2); \draw[g2p5-rt] (b2) -- (b3); \draw[g2p5-re] (b1) -- (b4);
        \node[font=\scriptsize\bfseries] at (1.4,1.5) {reverse search from 3};
    \end{scope}
    \node[note, text width=3.6cm, align=left, anchor=west] at (8.2,0.1) {template 3 prints\\\texttt{0: | 1: | 2: 0 1 |}\\\texttt{3: 0 1 2 | 4: 1 |}\\[2pt]one search per vertex; 4 is not an ancestor of 3 although both are below 1.};
\end{tikzpicture}
\figcap{Ancestors of 3 = the vertices marked by a search from 3 along reversed edges.}
\end{diagrambox}

In a DAG there is a faster-looking alternative: process the vertices in topological order (Chapter~5) and set $\mathrm{anc}(v) = \bigcup_{u \to v} \big(\mathrm{anc}(u) \cup \{u\}\big)$ with bitsets, \bigO{m \cdot n / 64}. For $n = 1000$, $m = 2000$ it ran in 7\,ms against 13\,ms for the $n$ reverse searches --- both instantaneous. The bitsets need $n^2/8$ bytes, though: fine for $n = 10^4$ (12.5\,MB), impossible for $n = 10^5$. And the per-vertex searches also work when the graph has cycles.

%------------------------------------------------------------
\subsubsection*{Reverse peeling: where can a walk go on forever?}

A vertex of out-degree 0 is a dead end: every walk that enters it stops. Delete it. Its \emph{predecessors} lose one exit each, and some of them may become dead ends in turn --- the reversed edges tell us exactly whose out-degree drops. When nothing more can be deleted, every remaining vertex has an exit to another remaining vertex, so a walk can go on forever; every deleted vertex leads only into dead ends. This is Kahn's algorithm (Chapter~5) run on $G^R$: peeling in-degree-0 vertices of $G^R$ $=$ peeling out-degree-0 vertices of $G$ (template 4, Exercise E5.4).

\begin{diagrambox}
\begin{tikzpicture}[x=1.25cm, y=1cm]
    \node[gvsel] (p1) at (0,0.8) {1}; \node[gvsel] (p2) at (1.2,1.4) {2}; \node[gvsel] (p3) at (1.2,0.1) {3};
    \node[gvdead] (p4) at (2.6,0.1) {4}; \node[gvdead] (p5) at (4.0,0.1) {5};
    \node[gvdead] (p6) at (2.6,-1.3) {6}; \node[gvdead] (p7) at (4.0,-1.3) {7};
    \draw[gde] (p1) -- (p2); \draw[gde] (p2) -- (p3); \draw[gde] (p3) -- (p1);
    \draw[gfaded, -{Stealth[length=2.2mm]}] (p3) -- (p4); \draw[gfaded, -{Stealth[length=2.2mm]}] (p4) -- (p5);
    \draw[gfaded, -{Stealth[length=2.2mm]}] (p6) -- (p4); \draw[gfaded, -{Stealth[length=2.2mm]}] (p6) -- (p7);
    \draw[gfaded, -{Stealth[length=2.2mm]}] (p7) -- (p5);
    \node[gdist, text=recursionred, above=1pt of p5] {round 1};
    \node[gdist, text=recursionred, above=1pt of p4] {round 2};
    \node[gdist, text=recursionred, right=1pt of p7] {round 2};
    \node[gdist, text=recursionred, left=1pt of p6] {round 3};
    \node[note, text width=5.4cm, align=left, anchor=west] at (5.2,0.1) {5 has out-degree 0. Removing it drops 4 and 7 to out-degree 0; removing those drops 6 to 0. Vertex 3 keeps its edge $3 \to 1$.\\[2pt]Template 4 deletes $5, 7, 4, 6$ and prints \texttt{1 2 3}: only the cycle and what still leads into it survive.};
\end{tikzpicture}
\figcap{Reverse peeling. Grey: deleted vertices; purple: vertices from which an endless walk exists.}
\end{diagrambox}

%------------------------------------------------------------
\subsubsection*{Variations map}

\begin{diagrambox}
\begin{tikzpicture}[
    var/.style={rectangle, rounded corners=2pt, draw=patternborder, fill=patternbg,
                font=\scriptsize, align=center, text width=2.95cm, minimum height=1.3cm, inner sep=3pt}]
    \node[var, fill=answerfill, draw=accentorange, text width=5.4cm, font=\scriptsize\bfseries] (root) at (0,0)
        {who can reach the target? search FROM the target in $G^R$};
    \node[var] (v1) at (-5.1,-2) {\textbf{all ancestors}\\one reverse search each\\(2.5.1)};
    \node[var] (v2) at (-1.7,-2) {\textbf{good in, bad out}\\wall the bad, search from the exit (2.5.2)};
    \node[var] (v3) at (1.7,-2)  {\textbf{from the border}\\enclosed $=$ unmarked\\(E5.1, E5.3)};
    \node[var] (v4) at (5.1,-2)  {\textbf{two targets}\\two searches, intersect\\(E5.2)};
    \node[var] (v5) at (-5.1,-3.8) {\textbf{endless walks}\\peel dead ends backwards\\(E5.4)};
    \node[var] (v6) at (-1.7,-3.8) {\textbf{a target set $T$}\\all of $T$ on the stack\\(E5.5)};
    \node[var] (v7) at (1.7,-3.8)  {\textbf{only cheap edges}\\search inside a binary search (E5.6)};
    \node[var] (v8) at (5.1,-3.8)  {\textbf{strongly connected?}\\one search in $G$, one in $G^R$ (template 5)};
    \node[draw=patternborder, dashed, rounded corners=4pt, fit=(v1)(v4)(v5)(v8), inner sep=5pt] (grp) {};
    \draw[flow] (root.south) -- (grp.north);
    \node[note, below=2mm of grp] {distances to the target: BFS / Dijkstra on $G^R$ (Chapters 3, 7); all pairs: SCCs + bitsets (Chapter 19)};
\end{tikzpicture}
\end{diagrambox}

\textbf{Twists and how to handle them}
\begin{itemize}[leftmargin=*]
    \item \textbf{The rule, not the graph, is directed} (``flow to a cell that is not higher''): reverse the rule (``climb to a cell that is not lower''), never build $G^R$ explicitly.
    \item \textbf{Targets are a whole region} (the border, all cells of one colour): push them all before the loop.
    \item \textbf{Two different targets} (two oceans, ``reaches $a$ and reaches $b$''): two reverse searches and an AND of the marks.
    \item \textbf{Bad starts must not escape}: neutralise them first (wall their neighbourhood), then search once from the target.
    \item \textbf{Output per vertex is a list}: one reverse search per vertex, \bigO{n(n+m)}; collect by scanning $u$ in order to get sorted lists without sorting.
    \item \textbf{``Forever'', ``safe'', ``doomed''}: peel out-degree-0 vertices using the reversed edges.
    \item \textbf{Strongly connected?} ``Everyone reaches 1 and 1 reaches everyone'' is one search from 1 in $G$ and one in $G^R$; both must mark all $n$ vertices (on $D$ the template marks 7 and 5: $D$ is not strongly connected).
\end{itemize}

%------------------------------------------------------------
\subsubsection*{Code templates}

\begin{lstlisting}[style=cpp]
#include <bits/stdc++.h>
using namespace std;

// (1) WHO CAN REACH T?  One iterative DFS in the reversed graph, from ALL targets at once.
//     radj[v] lists the u with an edge u -> v. Returns can[u] = 1 iff u reaches some target.
vector<char> canReach(const vector<vector<int>>& radj, const vector<int>& targets) {
    vector<char> can(radj.size(), 0);
    vector<int> st;
    for (int t : targets) if (!can[t]) { can[t] = 1; st.push_back(t); }
    while (!st.empty()) {
        int v = st.back(); st.pop_back();
        for (int u : radj[v])                    // reversed edge v -> u = original u -> v
            if (!can[u]) { can[u] = 1; st.push_back(u); }
    }
    return can;
}

// (2) ROLES REVERSED IN A GRID: mark every '.' cell connected to the border, from the border.
vector<string> markFromBorder(vector<string> g) {
    int R = g.size(), C = g[0].size();
    int dr[] = {1, -1, 0, 0}, dc[] = {0, 0, 1, -1};
    vector<pair<int, int>> st;
    for (int r = 0; r < R; r++)
        for (int c = 0; c < C; c++) {
            bool border = (r == 0 || c == 0 || r == R - 1 || c == C - 1);
            if (border && g[r][c] == '.') { g[r][c] = 'o'; st.push_back({r, c}); }
        }
    while (!st.empty()) {
        int r = st.back().first, c = st.back().second; st.pop_back();
        for (int d = 0; d < 4; d++) {
            int nr = r + dr[d], nc = c + dc[d];
            if (nr < 0 || nc < 0 || nr >= R || nc >= C || g[nr][nc] != '.') continue;
            g[nr][nc] = 'o'; st.push_back({nr, nc});
        }
    }
    return g;                                    // 'o' = reaches the border, '.' left = enclosed
}

// (3) ALL ANCESTORS: one reverse search per vertex, O(n (n + m)). Lists come out sorted.
vector<vector<int>> allAncestors(const vector<vector<int>>& radj) {
    int n = radj.size();
    vector<vector<int>> anc(n);
    for (int v = 0; v < n; v++) {
        vector<char> can = canReach(radj, vector<int>(1, v));
        for (int u = 0; u < n; u++) if (u != v && can[u]) anc[v].push_back(u);
    }
    return anc;
}

// (4) REVERSE PEELING: vertices from which an infinite walk exists.
//     Delete vertices of out-degree 0; the reversed edges tell whose out-degree drops.
vector<char> endless(const vector<vector<int>>& adj, const vector<vector<int>>& radj) {
    int n = adj.size();
    vector<int> out(n), q;
    vector<char> alive(n, 1);
    for (int v = 0; v < n; v++) { out[v] = adj[v].size(); if (!out[v]) q.push_back(v); }
    while (!q.empty()) {
        int v = q.back(); q.pop_back();
        alive[v] = 0;                            // v is a dead end: every walk from v stops
        for (int u : radj[v]) if (--out[u] == 0) q.push_back(u);
    }
    return alive;
}

int main() {
    // the graph D of Part A, 0-indexed (vertex k is k-1)
    vector<vector<int>> D = {{1, 4, 5}, {2, 3}, {0}, {2}, {3, 5}, {6}, {5}, {6}};
    int n = D.size();
    vector<vector<int>> R(n);
    for (int u = 0; u < n; u++) for (int v : D[u]) R[v].push_back(u);
    vector<char> a = canReach(R, {3});
    for (int u = 0; u < n; u++) if (a[u]) cout << u + 1 << " ";      // 1 2 3 4 5
    cout << "\n";
    vector<char> b = canReach(R, {1, 7});                              // targets 2 and 8
    for (int u = 0; u < n; u++) if (b[u]) cout << u + 1 << " ";      // 1 2 3 4 5 8
    cout << "\n";
    vector<string> g = markFromBorder({"#####", "#..#.", "#.#..", "#####"});
    for (auto& s : g) cout << s << " ";                               // ##### #..#o #.#oo #####
    cout << "\n";
    // a DAG: 0->2, 1->2, 2->3, 1->4
    vector<vector<int>> Rd = {{}, {}, {0, 1}, {2}, {1}};
    vector<vector<int>> an = allAncestors(Rd);
    for (int v = 0; v < 5; v++) {                // 0: | 1: | 2: 0 1 | 3: 0 1 2 | 4: 1 |
        cout << v << ":";
        for (int u : an[v]) cout << " " << u;
        cout << " | ";
    }
    cout << "\n";
    // P: a cycle 1 -> 2 -> 3 -> 1 with a tail 3 -> 4 -> 5 and 6 -> 4, 6 -> 7 -> 5
    vector<vector<int>> P = {{1}, {2}, {0, 3}, {4}, {}, {3, 6}, {4}}, PR(7);
    for (int u = 0; u < 7; u++) for (int v : P[u]) PR[v].push_back(u);
    vector<char> e = endless(P, PR);
    for (int u = 0; u < 7; u++) if (e[u]) cout << u + 1 << " ";      // 1 2 3
    cout << "\n";
    // (5) "1 reaches everyone and everyone reaches 1": one search in D, one in D^R
    vector<char> f = canReach(D, {0}), r = canReach(R, {0});
    cout << count(f.begin(), f.end(), 1) << " " << count(r.begin(), r.end(), 1) << "\n";  // 7 5
}
\end{lstlisting}

Note the last two lines: \texttt{canReach} is just a search, so called with \texttt{D} instead of \texttt{R} it is the ordinary forward search. From 1, $D$ reaches 7 vertices (not 8) and only 5 vertices reach 1 --- two reasons why $D$ is not strongly connected.

\begin{insightbox}[Which template when?]
\renewcommand{\arraystretch}{1.25}
\begin{tabularx}{\linewidth}{@{}>{\bfseries}l X@{}}
(1) \texttt{canReach} & ``who can reach $t$ / any vertex of $T$?''; with \texttt{adj} instead of \texttt{radj} it is the forward search (template 5 uses both) \\
(2) \texttt{markFromBorder} & grids: enclosed regions, lakes, ``surrounded''; with a height rule it becomes the ocean search of E5.2 \\
(3) \texttt{allAncestors} & a list per vertex, $n \le$ a few thousand; for DAGs with $n \le 10^4$ bitsets in topological order are an option \\
(4) \texttt{endless} & ``walk forever'', ``safe states'', ``which positions are doomed'' --- Kahn on $G^R$ \\
\end{tabularx}
\end{insightbox}

%------------------------------------------------------------
\subsubsection*{Common mistakes}

\begin{warningbox}[Mistake 1 --- building the reversed graph the wrong way round]
\begin{lstlisting}[style=snippet]
radj[a].push_back(b);     // WRONG: this is the original edge a -> b again
radj[b].push_back(a);     // RIGHT: edge a -> b is stored at b, pointing back to a
\end{lstlisting}
On $D$ the wrong version answers ``who can reach 4?'' with $\{3, 4, 1, 2, 5, 6, 7\}$ --- the vertices reachable \emph{from} 4.
\end{warningbox}

\begin{warningbox}[Mistake 2 --- reversing the rule only halfway in a grid]
``Water flows from $a$ to a neighbour $b$ with $h_b \le h_a$.'' The reverse search goes from $b$ to $a$ when $h_a \ge h_b$. Writing \texttt{h[nr][nc] <= h[r][c]} in the reverse search runs the \emph{forward} rule from the ocean: it marks the cells that water from the coast could flow \emph{to}, which is a different set.
\end{warningbox}

\begin{warningbox}[Mistake 3 --- one search per start ``because it is simpler'']
For $n = 10^5$, $m = 2 \cdot 10^5$ random edges and three targets, a forward search from every vertex (stopping as soon as a target is found) needed $3.8 \cdot 10^8$ stack pops and 63--95 seconds (three runs); the reverse search took 8--12\,ms (Exercise E5.5).
\end{warningbox}

\begin{warningbox}[Mistake 4 --- forgetting that the targets are marked too]
Template 1 marks the targets themselves. If the question is ``which \emph{other} vertices'' (ancestors, E5.6 ``every node except 0''), exclude them when collecting the answer, or count them correctly.
\end{warningbox}

\begin{warningbox}[Mistake 5 --- peeling with in-degrees of $G$]
For ``endless walks'' the dead ends are vertices of out-degree 0, and when one is deleted the out-degree of its \emph{predecessors} drops: iterate over \texttt{radj}. Peeling in-degree-0 vertices of $G$ answers a different question: it keeps the vertices that can be reached \emph{from} a cycle, not those that can reach one (on the graph of the figure it would keep $1, \dots, 5$ and delete 6 and 7).
\end{warningbox}

\begin{checklistbox}[``I'm stuck'' checklist for Pattern 5]
\begin{itemize}[leftmargin=*]
    \item[$\square$] Is the question asked from many starts towards few targets? Then search from the targets.
    \item[$\square$] Directed graph: did I store every edge $a \to b$ in \texttt{radj[b]}? Directed rule in a grid: did I flip the comparison?
    \item[$\square$] Are there several targets (a border, an ocean, a set)? Push them all before the loop.
    \item[$\square$] Two conditions (``reaches $a$ and $b$'')? Two reverse searches and an AND.
    \item[$\square$] Bad starts that must not escape? Neutralise them first, then search once.
    \item[$\square$] ``Forever'' or ``safe''? Peel out-degree-0 vertices along reversed edges.
    \item[$\square$] Do I need distances rather than reachability? Then it is BFS (Chapter~3), not DFS.
\end{itemize}
\end{checklistbox}
